\section{Type System}
\label{sec:grust:typing}

This section details the static type system of \GRust, which enforces the correct composition
of its synchronous and asynchronous constructs. A key design choice in \GRust is to implement a
dedicated type-checking pass rather than relying on the \Rust compiler. This approach is motivated
by several factors.

First, \GRust introduces domain-specific abstractions, such as \emph{signals} and \emph{events},
that have no equivalent in the \Rust type system. A custom type checker is essential to enforce
rules governing their use---for example, ensuring that an operator like \kft{timeout} receives an
event, not a signal.
%
Second, it enables the compiler to report clear, high-level errors that refer directly to the source
\GRust model. This shields the developer from cryptic error messages related to the underlying
generated code, which would otherwise report a \emph{mismatched types} error on unfamiliar code.
%
Finally, \GRust mandates explicit type annotations, deliberately omitting type inference. This
design promotes code clarity and predictability, which are paramount in safety-critical domains.
While this choice prevents the direct use of externally defined \Rust types, the benefits of
domain-specific validation and improved diagnostics are a worthwhile trade-off.

\begin{table}[!ht]
    \centering
    \begin{tabularx}{\textwidth}{|l|X|}\hline
        $\ID_s, \: \ID_f, \: \ID_c, \: \ID$     & Service, function, component,
        and variable identifiers respectively.                                           \\\hline
        $\IDs ::= \{\ID^*\}$                    & A set of identifiers.                  \\\hline
        $\G[\ID_f]$                             & A context mapping function
        identifiers to their definitions.                                                \\\hline
        $\G[\ID_c]$                             & A context mapping component
        identifiers to their definitions.                                                \\\hline
        $\TY   ::= \boolTY \: \vert \: \intTY
            \: \vert \: \floatTY \: \vert \:
        \unitTY$                                & Basic types for values.                \\\hline
        $\kind{\TY} ::= \signal{\TY} \: \vert
        \: \event{\TY}$                         & Flow types, representing a
        flow of values of type \TY.                                                      \\\hline
        $\STY ::= \TY \: \vert \: \opt{\TY}$    & Sampling types for flows:
        signals always have values, while events might be absent, making their
        values optional (with the `\opt{}').                                             \\\hline
        $\tyCtx[\ID]$                           & A typing context \tyCtx
        mapping identifiers to their types. It can be a \kind{\TY} for \GRFRP
        identifiers and \STY for \GRSync identifiers.                                    \\\hline
        \typing{elem}{\texttt{Elem}}            & Asserts that the \GRust
        construct \texttt{Elem} is well-typed with respect to \tyCtx.                    \\\hline
        $\typing{elem}{\texttt{Elem}}: [\STY^* \:
        \vert \: (\kind{\TY})^*]$               & Asserts that the \GRust
        construct \texttt{Elem} has the type $[\STY^*]$ or $[(\kind{\TY})^*]$ in \tyCtx. \\\hline
        $\CST : \TY$                            & The constant \CST has the
        type \TY.                                                                        \\\hline
        $\OP : [\TY_i^*] \rightarrow [\TY_o^*]$ & An operator \OP that takes
        inputs of types $[\TY_i^*]$ and returns outputs of types $[\TY_o^*]$.            \\\hline
        \defines{\Eq}{\IDs}                     & The set of identifiers \IDs
        is defined by the equation \Eq.                                                  \\\hline
    \end{tabularx}
    \caption[Notations for the type system of \GRust.]%
    {Notations for the logic of the type system of \GRust.}
    \label{tab:grust:typing:notations}
\end{table}

To formalize the type system, we introduce several notations, summarized in
\Cref{tab:grust:typing:notations}. At the core of the type system is the \emph{typing context},
denoted by \tyCtx, which acts as a symbol table that maps each identifier to its type. Because
\GRust integrates two distinct semantic models, the type system must handle two corresponding kinds
of types: \emph{flow types} for the asynchronous \GRFRP world and \emph{sampling types} for the
synchronous \GRSync world.

In \GRFRP services, identifiers correspond to dataflows that evolve over time. Their types, called
\emph{flow types} and denoted by \kind{\TY}, describe the temporal nature of the flow. A
\emph{signal} (\signal{\TY}) represents a value that is continuously available, while an
\emph{event} (\event{\TY}) represents a discrete occurrence. The underlying type of the values
carried by the flow is a \emph{basic type} (\TY), such as \boolTY, \intTY, or \floatTY.

In contrast, \GRSync components operate at discrete logical instants. When a component executes, it
computes on a \emph{sample} of its input flows. The types of these instantaneous values are called
\emph{sampling types}, denoted by \STY.
%
When a signal is sampled, it always yields a value, so its sampling type is its underlying basic
type \TY.
%
When an event is sampled, a value is present only if the event occurs at that specific instant. This
`presence or absence' is captured by an \emph{optional type}, \opt{\TY}. Thus, the sampling type of
an event is \opt{\TY}.

Finally, the formal rules use the judgment \typing{elem}{\text{Elem}} to assert that a \GRust
construct \texttt{Elem} is well-typed within a given context. The helper notation
\defines{\Eq}{\IDs} is used to retrieve the set of identifiers \IDs that are newly defined by an
equation \Eq.

\subsection{\GRust Programs}
\label{sec:grust:typing:prog}

A \GRust program is well-typed if all its constituent items are well-typed, as defined by the
judgment \progTyped{\Prog} in rule \nameref{rule:grust:typing:prog:program}.

\begin{center}
    \begin{minipage}{0.9\textwidth}
        \centering
        \begin{mathpar}
            %% Program
            \inferTY
            { \left(\itemTyped{\Item}\right)^* }
            { \progTyped{\Item^*} }
            {Prog}
            \label{rule:grust:typing:prog:program}
        \end{mathpar}
    \end{minipage}
\end{center}

An item is well-typed if it satisfies the judgment \itemTyped{\Item}. Imports and exports are
well-typed if the identifiers they define are mapped to flow types (\kind{\TY}) in the context
\tyCtx, as stated by rules \nameref{rule:grust:typing:prog:interface:import} and
\nameref{rule:grust:typing:prog:interface:export}.
%
Services, functions, and components are well-typed if their bodies are well-typed, as shown in rules
\nameref{rule:grust:typing:prog:service}, \nameref{rule:grust:typing:prog:fun}, and
\nameref{rule:grust:typing:prog:comp}. Additionally, the input and output types for functions and
components must be correctly recorded in the typing context. Functions are restricted to basic types
to forbid direct handling of events, which may require state.

\begin{center}
    \begin{minipage}{0.9\textwidth}
        \centering
        \begin{mathpar}
            %% Interface
            % Import
            \inferTY
            { \tyCtx[\ID] = \kind{\TY} }
            { \itemTyped{\importX{\typed{\FlowKind}{\ID}{\TY}}} }
            {Import}
            \label{rule:grust:typing:prog:interface:import}

            % Export
            \inferTY
            { \tyCtx[\ID] = \kind{\TY} }
            { \itemTyped{\exportX{\typed{\FlowKind}{\ID}{\TY}}} }
            {Export}
            \label{rule:grust:typing:prog:interface:export}

            %% Service
            \inferTY
            { \left(\stmtTyped{\FlowStmt}\right)^* }
            { \itemTyped{\servItem{\ID_s}{\interval{\INT}{\INT}}{\many{\FlowStmt}}} }
            {Serv}
            \label{rule:grust:typing:prog:service}

            %% Component
            \inferTY
            {
                \left(\tyCtx[\ID_i] = \STY_i\right)^* \\
                \eqTyped{\setEq} \\
                \left(\tyCtx[\ID_o] = \STY_o\right)^*
            }
            { \itemTyped{\compItem{\ID_c}{\many{(\ID_i \, \kft{:} \, \STY_i)}}{\many{(\ID_o \, \kft{:} \, \STY_o)}}{\many{\Eq}}} }
            {Comp}
            \label{rule:grust:typing:prog:comp}

            %% Function
            \inferTY
            {
                \left(\tyCtx[\ID_i] = \TY_i\right)^* \\
                \left(\stmtTyped{\Decl}\right)^* \\
                \exprTyped{\Expr}{\TY_o}
            }
            { \itemTyped{\funItem{\ID_f}{\many{(\ID_i \: \kft{:} \: \TY_i)}}{\TY_o}{\many{\Decl}}} }
            {Fun}
            \label{rule:grust:typing:prog:fun}
        \end{mathpar}
    \end{minipage}
\end{center}

\subsection{Bodies of \GRFRP Services}
\label{sec:grust:typing:grfrp}

This section presents the typing rules for services and other \GRFRP constructs.

\subsubsection{Flow Statements}
\label{sec:grust:typing:grfrp:flow_stmt}

A flow statement is well-typed if it satisfies the judgment \stmtTyped{\FlowStmt}.
%
A \kft{let} statement, following rule \nameref{rule:grust:typing:grfrp:flow_stmt:let}, is well-typed
if its flow operation (\FlowOp) produces types that match both the declared types ($[\kind{\TY}^*]$)
and the types stored in the context ($[(\tyCtx[\ID])^*]$).
%
Similarly, rule \nameref{rule:grust:typing:grfrp:flow_stmt:def} states that a definition statement
is well-typed if the types from its flow operation match those in the context.

\begin{center}
    \begin{minipage}{0.9\textwidth}
        \centering
        \begin{mathpar}
            %% Flow Statement
            % Let
            \inferTYFlowStmt
            {
                \flowTyped{\FlowOp}{(\tyCtx[\ID])^*} \\
                \left(\tyCtx[\ID] = \kind{\TY}\right)^*
            }
            { \stmtTyped{\letStmt{\many{(\typed{\FlowKind}{\ID}{\TY})}}{\FlowOp}} }
            {Let}
            \label{rule:grust:typing:grfrp:flow_stmt:let}

            % Definition
            \inferTYFlowStmt
            { \flowTyped{\FlowOp}{(\tyCtx[\ID])^*} }
            { \stmtTyped{\defStmt{\many{\ID}}{\FlowOp}} }
            {Def}
            \label{rule:grust:typing:grfrp:flow_stmt:def}
        \end{mathpar}
    \end{minipage}
\end{center}

\subsubsection{Flow Operations}
\label{sec:grust:typing:grfrp:flow_op}

A flow operation \FlowOp has the flow type \kind{\TY} if it satisfies the judgment
\flowTyped{\FlowOp}{\kind{\TY}}.
%
An identifier referring to a flow has the type stored in the context \tyCtx, as defined by the rule
\nameref{rule:grust:typing:grfrp:flow_op:ident}.
%
Function and component applications are well-typed if their input operations are compatible with the
signatures of the callees, as shown in rules \nameref{rule:grust:typing:grfrp:flow_op:funapp} and
\nameref{rule:grust:typing:grfrp:flow_op:compapp}.
%
Type compatibility is defined by rules \nameref{rule:grust:typing:grfrp:flow_op:sty:signal} and
\nameref{rule:grust:typing:grfrp:flow_op:sty:event}: a signal \signal{\TY} is compatible with its
inner type \TY, while an event \event{\TY} is compatible with an option of its inner type \opt{\TY}.

\begin{center}
    \begin{minipage}{0.99\textwidth}
        \centering
        \begin{mathpar}
            %% Flow Operation
            % Identifier
            \inferTYFlow
            { \tyCtx[\ID] = \kind{\TY} }
            { \flowTyped{\ID}{\kind{\TY}} }
            {ID}
            \label{rule:grust:typing:grfrp:flow_op:ident}

            % Function
            \inferTYFlow
            {
                \left(\flowTyped{\FlowOp}{\kind[i]{\TY_i}}\right)^* \\
                \G[\ID_f] = \funItem{\ID_f}{\many{(\ID_i \: \kft{:} \: \TY_i)}}{\TY_o}{\many{\Decl}} \\
                \left(\kind[i]{\TY_i} \doteq \TY_i\right)^* \\
                \kind[o]{\TY_o} \doteq \TY_o
            }
            { \flowTyped{\compApp{\ID_f}{\many{\FlowOp}}}{\kind[o]{\TY_o}} }
            {AppF}
            \label{rule:grust:typing:grfrp:flow_op:funapp}

            %% Application
            % Component
            \inferTYFlow
            {
                \left(\flowTyped{\FlowOp}{\kind[i]{\TY_i}}\right)^* \\
                \G[\ID_c] = \compItem{\ID_c}{\many{(\ID_i \, \kft{:} \, \STY_i)}}{\many{(\ID_o \, \kft{:} \, \STY_o)}}{\many{\Eq}} \\
                \left(\kind[i]{\TY_i} \doteq \STY_i\right)^* \\
                \left(\kind[o]{\TY_o} \doteq \STY_o\right)^*
            }
            { \flowTyped{\compApp{\ID_c}{\many{\FlowOp}}}{(\kind[o]{\TY_o})^*} }
            {AppC}
            \label{rule:grust:typing:grfrp:flow_op:compapp}
        \end{mathpar}
    \end{minipage}

    \begin{minipage}{0.9\textwidth}
        \centering
        \begin{mathpar}
            %% Type correspondence
            % Signal
            \inferTY
            { }
            {\signal{\TY} \doteq \TY}
            {Signal}
            \label{rule:grust:typing:grfrp:flow_op:sty:signal}

            % Event
            \inferTY
            { }
            {\event{\TY} \doteq \opt{\TY}}
            {Event}
            \label{rule:grust:typing:grfrp:flow_op:sty:event}
        \end{mathpar}
    \end{minipage}
\end{center}

The typing rules for the built-in \GReact operators enforce their correct usage with specific flow
types, as their semantics depends on the kind of flows. Each rule specifies the expected input and
output flow types, ensuring that operations like converting signals to events or
filtering data are performed safely.

The \kft{on\_change} and \kft{persist} operators bridge the gap
between signals and events. As shown in rule \nameref{rule:grust:typing:grfrp:flow_op:onchange},
\kft{on\_change} takes a \signal{\TY} and produces an \event{\TY} that fires only when the signal
value changes. Conversely, rule \nameref{rule:grust:typing:grfrp:flow_op:persist} shows that
\kft{persist} takes an \event{\TY} and creates a \signal{\TY} that holds the value of the last event
occurrence.
%
The \kft{scan} operator enables periodic observations over a signal.
Rule \nameref{rule:grust:typing:grfrp:flow_op:scan} states that it takes a \signal{\TY} and an
integer (the scanning period), producing a new \signal{\TY}. Its counterpart for events,
\kft{sample}, observes occurrences of an event periodically, for instance, to downsample a
high-frequency event source.
%
The \kft{timeout} operator, defined in rule
\nameref{rule:grust:typing:grfrp:flow_op:timeout}, expects an \event{\TY} and an integer duration;
if the event does not occur within that duration, it emits an occurrence of its output
\event{\unitTY}. The \kft{time} operator, following rule
\nameref{rule:grust:typing:grfrp:flow_op:time}, provides the current time as a \signal{\floatTY}.
%
The \kft{merge} operator, governed by rule
\nameref{rule:grust:typing:grfrp:flow_op:merge}, combines two event flows of the same type into a
single flow that fires when either input fires. The \kft{throttle} operator, following rule
\nameref{rule:grust:typing:grfrp:flow_op:throttle}, filters the input signal, emitting a new value
only when the change exceeds a specified threshold \DELTA (which is useful for ignoring sensor
noise).

\begin{center}
    \begin{minipage}{0.99\textwidth}
        \centering
        \begin{mathpar}
            %% GReact
            % OnChange
            \inferTYFlow
            { \flowTyped{\FlowOp}{\signal{\TY}} }
            { \flowTyped{\onchangeExpr{\FlowOp}}{\event{\TY}} }
            {OChg}
            \label{rule:grust:typing:grfrp:flow_op:onchange}

            % Persist
            \inferTYFlow
            { \flowTyped{\FlowOp}{\event{\TY}} }
            { \flowTyped{\persistExpr{\FlowOp}}{\signal{\TY}} }
            {Prst}
            \label{rule:grust:typing:grfrp:flow_op:persist}

            % Scan
            \inferTYFlow
            {
                \flowTyped{\FlowOp}{\signal{\TY}} \\
                \exprTyped{\INT}{\intTY}
            }
            { \flowTyped{\scanExpr{\FlowOp}{\INT}}{\signal{\TY}} }
            {Scan}
            \label{rule:grust:typing:grfrp:flow_op:scan}

            % Sample
            \inferTYFlow
            {
                \flowTyped{\FlowOp}{\event{\TY}} \\
                \exprTyped{\INT}{\intTY}
            }
            { \flowTyped{\scanExpr{\FlowOp}{\INT}}{\event{\TY}} }
            {Samp}
            \label{rule:grust:typing:grfrp:flow_op:sample}

        \end{mathpar}
    \end{minipage}

    \begin{minipage}{0.9\textwidth}
        \centering
        \begin{mathpar}
            % Timeout
            \inferTYFlow
            {
                \flowTyped{\FlowOp}{\event{\TY}} \\
                \exprTyped{\INT}{\intTY}
            }
            { \flowTyped{\timeoutExpr{\FlowOp}{\INT}}{\event{\unitTY}} }
            {TOut}
            \label{rule:grust:typing:grfrp:flow_op:timeout}

        \end{mathpar}
    \end{minipage}

    \begin{minipage}{0.9\textwidth}
        \centering
        \begin{mathpar}
            % Merge
            \inferTYFlow
            {
                \flowTyped{\FlowOp_1}{\event{\TY}} \\
                \flowTyped{\FlowOp_2}{\event{\TY}}
            }
            { \flowTyped{\mergeExpr{\FlowOp_1}{\FlowOp_2}}{\event{\TY}} }
            {Mrg}
            \label{rule:grust:typing:grfrp:flow_op:merge}
        \end{mathpar}
    \end{minipage}

    \begin{minipage}{0.9\textwidth}
        \centering
        \begin{mathpar}
            % Throttle
            \inferTYFlow
            {
                \flowTyped{\FlowOp}{\signal{\TY}} \\
                \exprTyped{\DELTA}{\TY}
            }
            { \flowTyped{\throttleExpr{\FlowOp}{\DELTA}}{\signal{\TY}} }
            {Thrt}
            \label{rule:grust:typing:grfrp:flow_op:throttle}

            % Time
            \inferTYFlow
            { }
            { \flowTyped{\timeExpr}{\signal{\floatTY}} }
            {Time}
            \label{rule:grust:typing:grfrp:flow_op:time}
        \end{mathpar}
    \end{minipage}
\end{center}

\subsection{Bodies of \GRSync Functions and Components}
\label{sec:grust:typing:grsync}

This section describes the typing rules for the bodies of functions and components, and other
\GRSync constructs.

\subsubsection{Declarations and Equations}
\label{sec:grust:typing:grsync:eq}

An equation is well-typed if it satisfies the judgment \eqTyped{\Eq}. A set of equations is
well-typed if each individual equation is well-typed, as stated by rule
\nameref{rule:grust:typing:grsync:eqs}.
%
Declarations and definitions, following rules \nameref{rule:grust:typing:grsync:eq:decl} and
\nameref{rule:grust:typing:grsync:eq:def}, are well-typed if the type of the expression matches the
type of the identifier in the context.
%
An initialization is well-typed if the constant type matches the types of both the identifier and
its corresponding memory, as defined by rule \nameref{rule:grust:typing:grsync:eq:init}.

\begin{center}
    \begin{minipage}{0.9\textwidth}
        \centering
        \begin{mathpar}
            %% Set of equations
            \inferTYEq
            { \left(\eqTyped{\Eq}\right)^* }
            { \eqTyped{\setEq} }
            {Blck}
            \label{rule:grust:typing:grsync:eqs}

            %% Declarations
            \inferTYEq
            {
                \exprTyped{\Expr}{(\tyCtx[\ID])^*} \\
                \left(\tyCtx[\ID] = \STY\right)^*
            }
            { \eqTyped{\declEq{\many{(\ID \: \kft{:} \: \STY)}}{\Expr}} }
            {Decl}
            \label{rule:grust:typing:grsync:eq:decl}

            %% Definition
            \inferTYEq
            { \exprTyped{\Expr}{(\tyCtx[\ID])^*} }
            { \eqTyped{\defEq{\many{\ID}}{\Expr}} }
            {Def}
            \label{rule:grust:typing:grsync:eq:def}

            %% Initialization
            \inferTYEq
            {
                \exprTyped{\CST}{\tyCtx[\lastExpr{\ID}]} \\
                \tyCtx[\ID] = \tyCtx[\lastExpr{\ID}]
            }
            { \eqTyped{\initEq{\ID}{\CST}} }
            {Init}
            \label{rule:grust:typing:grsync:eq:init}
        \end{mathpar}
    \end{minipage}
\end{center}

The typing of a \kft{match} equation followes rule \nameref{rule:grust:typing:grsync:eq:match}. It
requires that for each branch, the pattern types match the expression types, the guard is boolean,
and the branch equations are well-typed.

\begin{center}
    \begin{minipage}{0.99\textwidth}
        \centering
        \begin{mathpar}
            %% Match
            \inferTYEq
            {
                \exprTyped{\Expr}{\STY^*} \\
                \left(\patTyped{\Pat_k}{\STY^*}\right)^* \\\\
                \left(\exprTyped{\Expr_k}{\boolTY}\right)^* \\
                \left(\eqTyped{\setEq[k]}\right)^*
            }
            { \eqTyped{\matchX{\Expr}{\many{(\normBr{\gPat{\Pat_k}{\Expr_k}}{\setEq[k]})}}} }
            {Mtch}
            \label{rule:grust:typing:grsync:eq:match}
        \end{mathpar}
    \end{minipage}
\end{center}

Similarly, a \kft{when} equation, following rule \nameref{rule:grust:typing:grsync:eq:when}, is
well-typed if its initialization constants match the signal types, and in each branch, the event
pattern is well-typed, the guard is boolean, and the equations are well-typed. Furthermore, any
event defined within the \kft{when} block must have an optional type.

\begin{center}
    \begin{minipage}{0.99\textwidth}
        \centering
        \begin{mathpar}
            %% When
            \inferTYEq
            {
                \left(\exprTyped{\CST}{\tyCtx[\lastExpr{\ID}]}\right)^* \\
                \left(\tyCtx[\ID] = \tyCtx[\lastExpr{\ID}]\right)^* \\\\
                \left(\epatTyped{\EPat_k}\right)^* \\
                \left(\exprTyped{\Expr_k}{\boolTY}\right)^* \\
                \eqTyped{\setEq[k]} \\\\
                \defines{\kft{when} \: \block{\dots}}{\IDs} \\
                \forall \ID_e \in \IDs \backslash \{\ID^*\}, \: \tyCtx[\ID_e] = \opt{\TY}
            }
            { \eqTyped{\whenEq{\initBr{\many{(\ID \: \kfm{=} \: \CST\kft{;})}}}{\many{(\normBr{\gPat{\EPat_k}{\Expr_k}}{\setEq[k]})}}} }
            {When}
            \label{rule:grust:typing:grsync:eq:when}
        \end{mathpar}
    \end{minipage}
\end{center}

\subsubsection{Expressions}
\label{sec:grust:typing:grsync:expr}

An expression \Expr has a sampling type $[\STY^*]$ if it satisfies the judgment
\exprTyped{\Expr}{\STY^*}.
%
Constants, identifiers, and \kft{last} memory accesses have the types assigned in the context, as
defined by rules \nameref{rule:grust:typing:grsync:expr:const},
\nameref{rule:grust:typing:grsync:expr:id}, and \nameref{rule:grust:typing:grsync:expr:last}.
%
An event emission \kft{emit}, following rule \nameref{rule:grust:typing:grsync:expr:emit}, wraps a
basic type \TY into an optional type \opt{\TY}.

\begin{center}
    \begin{minipage}{0.9\textwidth}
        \centering
        \begin{mathpar}
            %% Constant
            \inferTYE
            { \CST : \TY }
            { \exprTyped{\CST}{\TY} }
            {Cst}
            \label{rule:grust:typing:grsync:expr:const}

            %% Identifier
            \inferTYE
            { \tyCtx[\ID] = \STY }
            { \exprTyped{\ID}{\STY} }
            {ID}
            \label{rule:grust:typing:grsync:expr:id}

            %% Last
            \inferTYE
            { \tyCtx[\lastExpr{\ID}] = \STY }
            { \exprTyped{\lastExpr{\ID}}{\STY} }
            {Last}
            \label{rule:grust:typing:grsync:expr:last}

            %% Emit
            \inferTYE
            { \exprTyped{\Expr}{\TY} }
            { \exprTyped{\emitExpr{\Expr}}{\opt{\TY}} }
            {Emit}
            \label{rule:grust:typing:grsync:expr:emit}
        \end{mathpar}
    \end{minipage}
\end{center}

A \kft{match} expression followes rule \nameref{rule:grust:typing:grsync:expr:match}. It is
well-typed if the matched expression and patterns are type-compatible, guards are boolean, and all
branches yield expressions of the same type.

\begin{center}
    \begin{minipage}{0.9\textwidth}
        \centering
        \begin{mathpar}
            %% Match
            \inferTYE
            {
                \exprTyped{\Expr}{\STY_e^*} \\
                \left(\patTyped{\Pat_k}{\STY_e^*}\right)^* \\
                \left(\exprTyped{\Expr_k^g}{\boolTY}\right)^* \\
                \left(\exprTyped{\Expr_k}{\STY}\right)^*
            }
            { \exprTyped{\matchX{\Expr}{\many{(\normBr{\gPat{\Pat_k}{\Expr_k^g}}{\Expr_k})}}}{\STY} }
            {Mtch}
            \label{rule:grust:typing:grsync:expr:match}
        \end{mathpar}
    \end{minipage}
\end{center}

Applications of operators, functions, or components are well-typed if the input expression types and
output types align with the signatures of the callees, as stated by rules
\nameref{rule:grust:typing:grsync:expr:op}, \nameref{rule:grust:typing:grsync:expr:compapp}, and
\nameref{rule:grust:typing:grsync:expr:funapp}.

\begin{center}
    \begin{minipage}{0.99\textwidth}
        \centering
        \begin{mathpar}
            %% $n$-ary operation
            \inferTYE
            {
            \left(\exprTyped{\Expr}{\TY_i}\right)^* \\
            \OP : [\TY_i^*] \rightarrow [\TY_o^*]
            }
            { \exprTyped{\opExpr{\many{\Expr}}}{\TY_o^*} }
            {Op}
            \label{rule:grust:typing:grsync:expr:op}

            %% Application
            % Component
            \inferTYE
            {
                \left(\exprTyped{\Expr}{\STY_i}\right)^* \\
                \G[\ID_c] = \compItem{\ID_c}{\many{(\ID_i \, \kft{:} \, \STY_i)}}{\many{(\ID_o \, \kft{:} \, \STY_o)}}{\many{\Eq}}
            }
            { \exprTyped{\appExpr{\ID_c}{\many{\Expr}}}{(\STY_o)^*} }
            {AppC}
            \label{rule:grust:typing:grsync:expr:compapp}

            % Function
            \inferTYE
            {
                \left(\exprTyped{\Expr}{\TY_i}\right)^* \\
                \G[\ID_f] = \funItem{\ID_f}{\many{(\ID_i \: \kft{:} \: \TY_i)}}{\TY_o}{\many{\Decl}}
            }
            { \exprTyped{\appExpr{\ID_f}{\many{\Expr}}}{\TY_o} }
            {AppF}
            \label{rule:grust:typing:grsync:expr:funapp}
        \end{mathpar}
    \end{minipage}
\end{center}

\subsubsection{Patterns}
\label{sec:grust:typing:grsync:pat}

A pattern \Pat matches the types $[\STY^*]$ if it satisfies the judgment
\patTyped{\Pat}{\STY^*}. Constant and identifier patterns match the types of their corresponding
constants and identifiers, as defined by rules \nameref{rule:grust:typing:grsync:pat:const} and
\nameref{rule:grust:typing:grsync:pat:id}.

\begin{center}
    \begin{minipage}{0.9\textwidth}
        \centering
        \begin{mathpar}
            %% Constant
            \inferTYPat
            { \tyCtx[\CST] = \TY }
            { \patTyped{\CST}{\TY} }
            {Cst}
            \label{rule:grust:typing:grsync:pat:const}

            %% Identifier
            \inferTYPat
            { \tyCtx[\ID] = \STY }
            { \patTyped{\ID}{\STY} }
            {ID}
            \label{rule:grust:typing:grsync:pat:id}
        \end{mathpar}
    \end{minipage}
\end{center}

A wildcard pattern matches any type, and a tuple pattern is well-typed if each of its sub-patterns
is well-typed against the corresponding type in the tuple, following rules
\nameref{rule:grust:typing:grsync:pat:wild} and \nameref{rule:grust:typing:grsync:pat:tuple}.

\begin{center}
    \begin{minipage}{0.9\textwidth}
        \centering
        \begin{mathpar}
            %% Wildcard
            \inferTYPat
            { }
            { \patTyped{\_}{\STY} }
            {Wild}
            \label{rule:grust:typing:grsync:pat:wild}

            %% Tuple
            \inferTYPat
            { \left(\patTyped{\Pat}{\STY}\right)^* }
            { \patTyped{\tuple{\Pat}}{\STY^*} }
            {Tupl}
            \label{rule:grust:typing:grsync:pat:tuple}
        \end{mathpar}
    \end{minipage}
\end{center}

\subsubsection{Event Patterns}
\label{sec:grust:typing:grsync:epat}

An event pattern \EPat is well-typed if it satisfies the judgment \epatTyped{\EPat}.
%
The rule \nameref{rule:grust:typing:grsync:epat:event} ensures that in an event detection, the identifier
capturing the event value has a type matching the inner type of the optional event sample.
%
A rising edge pattern follows rule \nameref{rule:grust:typing:grsync:epat:edge} and requires a
boolean expression. A tuple of event patterns is well-typed if each sub-pattern is well-typed, as
defined by rule \nameref{rule:grust:typing:grsync:epat:tuple}.

\begin{center}
    \begin{minipage}{0.9\textwidth}
        \centering
        \begin{mathpar}
            %% Event detection
            \inferTYEPat
            {
                \tyCtx[\ID_e] = \opt{\TY} \\
                \tyCtx[\ID] = \TY
            }
            { \epatTyped{\ID \: \kfm{=} \: \opt{\ID_e}} }
            {Dtct}
            \label{rule:grust:typing:grsync:epat:event}

            %% Rising edge
            \inferTYEPat
            { \exprTyped{\Expr}{\boolTY} }
            { \epatTyped{\Expr} }
            {Edge}
            \label{rule:grust:typing:grsync:epat:edge}

            %% Tuple
            \inferTYEPat
            { \left(\epatTyped{\EPat}\right)^* }
            { \epatTyped{\tuple{\EPat}} }
            {Tupl}
            \label{rule:grust:typing:grsync:epat:tuple}
        \end{mathpar}
    \end{minipage}
\end{center}

\subsection{Summary}
This section has defined the formal type system of \GRust, which provides the static guarantees for
bridging the asynchronous \GRFRP kernel with the synchronous \GRSync kernel. The system enforces a
crucial distinction between the \emph{flow types} of services (\eg \signal{\TY} and \event{\TY})
and the \emph{sampling types} used within components (\eg \TY and \opt{\TY}). This discipline is
fundamental to writing correct and reliable reactive programs, preserving the distinct semantics
of signals (continuously available values) and events (discrete occurrences).

A key principle of the type system is the requirement for explicit type annotations. Unlike the
standard \Rust compiler, \GRust does not perform type inference. This design choice ensures that
type errors are reported at the level of the \GRust model, providing clear, high-level feedback that
is essential in a safety-critical context.

As future work, a formal proof of soundness would be valuable. Such a proof would demonstrate that
any well-typed \GRust program generates \Rust code that is guaranteed to be accepted by the \Rust
type checker, thus formally connecting the high-level model to the generated implementation.
